\section{Rotational covariance}\label{sec:rotations}

The rotation argument starts with a field insertion separated from
the two observables in a reflection form. Gaussian domination bounds
the insertion in the reconstructed Hilbert space, independently of
the number of factors in either observable. The spectral cone then
allows complex translations. Together these bounds give bounded
operators for insertions concentrated at a point of the
$(x_1,x_2)$-plane. Closedness removes the extra angular frequencies
at two points. At higher orders, reflection positivity gives a positive
kernel for the highest angular coefficient. Polygon divisibility and a
convexity estimate then reduce every order to the two-point result.


We use the Gaussian moment bounds, reconstruction and spectral cone
recalled in \cref{sec:probability}. The parent law and its components
have these properties by \cref{prop:limit-nine-mirrors}. Their
closedness enters the two-point frequency calculation and the
removal of contact terms.
Let $\mu$ be a translation-invariant, Gaussian-dominated law,
invariant and reflection positive at $x_1=0$.
We first work with this reconstruction, writing
$\mathcal H=\mathcal H_1$, $H=H_1$, and
$U_t=U(te_1)$.
Let $\mathcal A_+$ denote the positive-side measurable functions
belonging to every $L^p(\mu)$, $p<\infty$. It is an algebra.
Their classes are dense in $\mathcal H$, since $\mathcal A_+$
contains the bounded positive-side functions.

Use the Hodge isometry
$*f=(f_{23},-f_{13},f_{12})$ to write
$\mathbb B(b)=F(*^{-1}b)$, complex linearly. The completion recalled
in \cref{sec:probability} extends this map to vector-valued $L^2$ tests.
The action on these tests is
\[
 (g\cdot b)(x)=\det(g)\,g\,b(g^{-1}x)
 \quad(g\in O(3)),\qquad
 (\tau_a b)(x)=b(x-a).
\]
In particular
$\Theta_1\mathbb B(b)=\mathbb B(\overline{r_1\cdot b})$
and $\Theta_1U_t=U_{-t}\Theta_1$.
Write $L^2_{[a,b]}$ for the vector tests vanishing almost everywhere
outside $a\le x_1\le b$.

\subsection{The insertion estimate}\label{subsec:rotation-insertion}

For $\ell>0$ and $h\in L^2_{[0,\ell]}$, consider
$M_\ell(h)\Phi=e^{\mathbb B(h)}U_\ell\Phi$ for bounded
positive-side $\Phi$. This function lies in $\mathcal A_+$.
We first prove
\begin{equation}\label{eq:rotation-source-bound}
 \|[M_\ell(h)\Phi]_1\|
       \le e^{\|\Re h\|_2^2/2}\|[\Phi]_1\|.
\end{equation}

Set $h^\#=\tau_{\ell e_1}(\overline{r_1\cdot h}+h)$.
Reflection and translation invariance give
\begin{equation}\label{eq:rotation-source-doubling}
 \|[M_\ell(h)\Phi]_1\|^2
   =\langle[\Phi]_1,[M_{2\ell}(h^\#)\Phi]_1\rangle.
\end{equation}
The two summands in $h^\#$ have disjoint slab interiors, and
$\tau_{2\ell e_1}\overline{r_1\cdot h^\#}=h^\#$.

For a source $h'$ on $[0,\ell']$ with
$\tau_{\ell'e_1}\overline{r_1\cdot h'}=h'$, let
$h'^{(n)}=\sum_{j=0}^{n-1}\tau_{j\ell'e_1}h'$.
Doubling $h'^{(n)}$ gives $h'^{(2n)}$. Thus, with
$m_a=\|[M_{2^a\ell'}(h'^{(2^a)})\Phi]_1\|$,
\[
 m_a^2\le\|[\Phi]_1\|m_{a+1},\qquad
 m_0\le\|[\Phi]_1\|^{1-2^{-a}}m_a^{2^{-a}}.
\]
On the other hand, the integrand in the doubling identity is bounded
in modulus by
\[
 \|\Phi\|_\infty^2e^{\mathbb B(\Re h'^{(2^{a+1})})}.
\]
The real translates have disjoint interiors, so Gaussian domination
gives
\[
 m_a^2\le\|\Phi\|_\infty^2
       e^{\,2^a\|\Re h'\|_2^2}.
\]
Taking the indicated powers and sending $a\to\infty$ proves
\eqref{eq:rotation-source-bound} for symmetric sources, including
$\|[\Phi]_1\|=0$. Apply this to $h^\#$, whose real part has squared
norm $2\|\Re h\|_2^2$, in
\eqref{eq:rotation-source-doubling}. Cauchy--Schwarz proves the
bound for arbitrary $h$.

\begin{proposition}\label{prop:rotation-insertion}
For $\ell>0$ and $f\in L^2_{[0,\ell]}$ there is a unique bounded
operator $\mathcal B_\ell(f)$ on $\mathcal H$ satisfying
\[
 \mathcal B_\ell(f)[\Phi]_1
       =[\mathbb B(f)U_\ell\Phi]_1\quad(\Phi\in\mathcal A_+).
\]
It is complex linear in $f$, and
\begin{equation}\label{eq:rotation-insertion-bound}
 \|\mathcal B_\ell(f)\|\le k_1\|f\|_2,\qquad
 k_1=\sqrt{2\exp(1)}.
\end{equation}
For real $f$, the constant is $\exp(1/2)$.
If $f\in L^2_{[c,c']}$ with $0\le c<c'\le\ell$, then
\begin{equation}\label{eq:rotation-insertion-local}
 \mathcal B_\ell(f)=e^{-cH}
       \mathcal B_{c'-c}(\tau_{-ce_1}f)e^{-(\ell-c')H}.
\end{equation}
For $f\in L^2_{[-d,d]}$, set
$\mathcal B_d^{\mathrm c}(f)=\mathcal B_{2d}(\tau_{de_1}f)$.
Its matrix elements and adjoint are
\begin{align}
 \langle[\Phi']_1,\mathcal B_d^{\mathrm c}(f)[\Phi]_1\rangle
   &=\E_\mu[(\Theta_1U_d\Phi')\mathbb B(f)U_d\Phi],
                         \label{eq:rotation-insertion-centered}\\
 \mathcal B_d^{\mathrm c}(f)^*
   &=\mathcal B_d^{\mathrm c}(\overline{r_1\cdot f}).\notag
\end{align}
For $r\ge|t|$ and $a\in\R$, one also has
\begin{multline}\label{eq:rotation-insertion-shift}
 e^{-(r+t)H+ia\mathsf P_2}\mathcal B_d^{\mathrm c}(f)
           e^{-(r-t)H-ia\mathsf P_2}\\
       =\mathcal B_{d+r}^{\mathrm c}(\tau_{te_1+ae_2}f).
\end{multline}
\end{proposition}

\begin{proof}
For real $f$ and bounded $\Phi$, the map
$z\mapsto[e^{z\mathbb B(f)}U_\ell\Phi]_1$ is entire.
Indeed its exponential series converges locally uniformly in
$L^2(\mu)$, dominated on $|z|\le R$ by
$\|\Phi\|_\infty e^{R|\mathbb B(f)|}$.
The quotient map is a contraction. On $|z|=\|f\|_2^{-1}$,
\eqref{eq:rotation-source-bound} bounds its norm by
$\exp(1/2)\|[\Phi]_1\|$. Cauchy's coefficient estimate gives the
real insertion bound. Applying it to $\Re f$ and $\Im f$, and using
$\|\Re f\|_2+\|\Im f\|_2\le\sqrt2\|f\|_2$, gives
\eqref{eq:rotation-insertion-bound}.
The estimate sends null classes to null classes, so it defines the
operator on the dense classes of bounded functions and extends it
uniquely to $\mathcal H$.

For $\Phi\in\mathcal A_+$, its bounded truncations converge in $L^4$.
Since $\mathbb B(f)\in L^4$, their products with the inserted field
converge in $L^2$ by H\"older's inequality. This proves the asserted
formula on $\mathcal A_+$ and complex linearity in $f$.

Equation \eqref{eq:rotation-insertion-local} follows by applying
$U_c$ to
$\mathbb B(\tau_{-ce_1}f)U_{\ell-c}\Phi$ and using
\exteq{6}{ErgodicComponents}. Translation by $-de_1$ in the reflection
form gives \eqref{eq:rotation-insertion-centered}. Applying reflection
and conjugation to that integrand gives the adjoint formula.
For \eqref{eq:rotation-insertion-shift}, apply the two exterior
semigroups to a class $[\Phi]_1$. The resulting function is
\[
 \mathbb B(\tau_{(d+r)e_1+te_1+ae_2}f)\,
       U_{2(d+r)}\Phi.
\]
The test $\tau_{te_1+ae_2}f$ is supported in $[-d-r,d+r]$,
since $|t|\le r$. The displayed function therefore defines its centered
insertion at $d+r$, which is the right side. Equality on these dense
classes proves the operator identity.
\end{proof}

In particular, for $\Phi_1,\Phi_2\in\mathcal A_+$ measurable on
$x_1>d$,
\begin{equation}\label{eq:rotation-separated-insertion}
 |\E_\mu[(\Theta_1\Phi_1)\mathbb B(f)\Phi_2]|
 \le k_1\|f\|_2\,\|[U_{-d}\Phi_1]_1\|
                         \|[U_{-d}\Phi_2]_1\|.
\end{equation}
This is \eqref{eq:rotation-insertion-centered} applied to the two
translated observables. The constant does not depend on their degrees.

\subsection{Complex translations and point plates}
\label{subsec:rotation-point-plates}

From now on assume also the sharp cone
$\mathsf E_1\{(E,p_\perp):E<|p_2|\}=0$.
For the laws under construction, Proposition~\extref{3.1}{ErgodicComponents} supplies this from the diagonal mirror.

For $\xi=(\xi^0,\xi^1)\in\C^2$, put
\[
 d_{\mathcal W}(\xi)=\Re\xi^0-|\Im\xi^1|,\qquad
 \mathcal W=\{d_{\mathcal W}>0\},\qquad
 \xi^*=(\overline{\xi^0},-\overline{\xi^1}),\quad n_0=(1,0).
\]
The cone \exteq{8}{ErgodicComponents} gives contractions
\[
 T(\xi)=\int e^{-\xi^0E+i\xi^1p_2}\,\dd\mathsf E_1(E,p_\perp)
                  \quad(\xi\in\overline{\mathcal W}).
\]
They obey $T(\xi)T(\xi')=T(\xi+\xi')$,
$T(\xi)^*=T(\xi^*)$, and $T(\xi)\Omega_1=\Omega_1$.
They are strongly continuous on the closed tube by spectral
dominated convergence. In its interior they are holomorphic in
operator norm: on $d_{\mathcal W}\ge\epsilon>0$ the first and
second derivatives of the multiplier are bounded by constants
times $E e^{-\epsilon E}$ and $E^2e^{-\epsilon E}$.
Taylor's formula therefore gives a uniform quadratic remainder.
For real $y=(y^0,y^1)$ with $y^0\ge0$,
$T(y)[\Phi]_1=[U(y^0e_1+y^1e_2)\Phi]_1$.

Write $x=(y,x_3)$ with $y=(x_1,x_2)$. Set
\begin{gather*}
 e^+=e_1+ie_2,\quad e^-=e_1-ie_2,\quad e^3=e_3,\\
 w(+)=1,\quad w(-)=-1,\quad w(3)=0,\\
 s(+)=s(-)=1,\quad s(3)=-1.
\end{gather*}
A label $\lambda=(\sigma,\chi)$ has
$\sigma\in\{+,-,3\}$ and real $\chi\in L^2(\R)$.
For $\alpha\in L^2(\R^2;\C)$ its plate is
$(\alpha\otimes\lambda)(x)=\alpha(y)\chi(x_3)e^\sigma$.
Thus $\|\alpha\otimes\lambda\|_2\le\sqrt2\|\alpha\|_2\|\chi\|_2$.
For rotations $R_t$ in the $(x_1,x_2)$-plane and the coordinate
reflection $r_1$, direct multiplication gives
\begin{align}
 R_t\cdot(\alpha\otimes\lambda)
   &=e^{-iw(\sigma)t}(\alpha\circ R_{-t})\otimes\lambda,
                              \label{eq:rotation-plate-actions}\\
 \overline{r_1\cdot(\alpha\otimes\lambda)}
   &=s(\sigma)(\overline\alpha\circ r_1)\otimes\lambda.\notag
\end{align}

\begin{proposition}\label{prop:rotation-point-plates}
For each label $\lambda$ there are bounded operators
$\mathcal B_\lambda(\xi,\xi')$, $\xi,\xi'\in\mathcal W$,
holomorphic in operator norm and real linear in $\chi$, such that
\begin{equation}\label{eq:rotation-point-bound}
 \|\mathcal B_\lambda(\xi,\xi')\|
 \le\frac{k_2\|\chi\|_2}
          {\min(d_{\mathcal W}(\xi),d_{\mathcal W}(\xi'))},
 \qquad k_2=\frac8\pi\exp(1/2).
\end{equation}
For $\kappa,\kappa'\in\overline{\mathcal W}$,
\begin{align}
 T(\kappa)\mathcal B_\lambda(\xi,\xi')T(\kappa')
   &=\mathcal B_\lambda(\kappa+\xi,\xi'+\kappa'),
                                   \label{eq:rotation-point-shift}\\
 \mathcal B_\lambda(\xi,\xi')^*
   &=s(\sigma)\mathcal B_\lambda(\xi'^*,\xi^*).\notag
\end{align}
If $\alpha\in L^2(\R^2)$ is compactly supported in $|y^0|<d$, then
\begin{equation}\label{eq:rotation-point-integral}
 \mathcal B_d^{\mathrm c}(\alpha\otimes\lambda)
   =\int\alpha(y)\mathcal B_\lambda(dn_0+y,dn_0-y)\,\dd y.
\end{equation}
The integral is a Bochner integral in operator norm.
When $\supp\alpha\subset\{|y^0|\le d'\}$, $d'<d$, its norm is at
most $k_2\|\chi\|_2\|\alpha\|_1/(d-d')$.
\end{proposition}

\begin{proof}
Let $j_\delta=(\pi\delta^2)^{-1}\mathbf1_{\{|y|<\delta\}}$.
Set $A_\delta=\mathcal B_\delta^{\mathrm c}(j_\delta\otimes\lambda)$
and, when both tube distances exceed $\delta$, set
\[
 \mathcal B_\lambda^\delta(\xi,\xi')
   =T(\xi-\delta n_0)A_\delta T(\xi'-\delta n_0).
\]
These operators obey the shift and adjoint identities in
\eqref{eq:rotation-point-shift}. Here $j_\delta$ is real and
reflection invariant, so $A_\delta^*=s(\sigma)A_\delta$.
The real translation identity \eqref{eq:rotation-insertion-shift}
also gives, for $r\ge|y^0|$,
\begin{equation}\label{eq:rotation-point-regularized-translation}
 \mathcal B_\lambda^\delta((r+\delta)n_0+y,(r+\delta)n_0-y)
    =\mathcal B_{\delta+r}^{\mathrm c}((\tau_yj_\delta)\otimes\lambda).
\end{equation}

We prove a bound uniform as $\delta\downarrow0$.
Fix $\xi,\xi'$ and set
$d_0=\min(d_{\mathcal W}(\xi),d_{\mathcal W}(\xi'))$,
$r=(d_0-\delta)/2>0$. For unit vectors $\psi,\psi'$,
\[
 h(z)=\langle\psi,\mathcal B_\lambda^\delta(\xi+z,\xi'-z)\psi'\rangle
\]
is holomorphic on $|\Re z^0|<r$, $|\Im z^1|<r$.
Write $z=y+i(t,t')$. Factoring the exterior translations gives
\[
 h(y+i(t,t'))=
 \langle\psi_1,\mathcal B_{\delta+r}^{\mathrm c}
             ((\tau_yj_\delta)\otimes\lambda)\psi'_1\rangle,
\]
where
$\psi_1=T(\xi-(\delta+r)n_0+i(t,t'))^*\psi$ and
$\psi'_1=T(\xi'-(\delta+r)n_0-i(t,t'))\psi'$.
For $|t'|<r$ these are contractions of the original vectors;
they are independent of $y$.
For $\eta\in L^2((-r,r)^2)$, extended by zero, the integral of
$\eta(y)\tau_yj_\delta$ is $\eta*j_\delta$ in $L^2$.
Its support lies in $|y^0|\le r+\delta$.
The insertion bound and Young's inequality therefore give
\[
 \left|\int\eta(y)h(y+i(t,t'))\,\dd y\right|
   \le2\exp(1/2)\|\eta\|_2\|\chi\|_2.
\]
Hence
$\int_{(-r,r)^2}|h(y+i(t,t'))|^2\,\dd y
\le4\exp(1)\|\chi\|_2^2$.
Apply the mean-value inequality on a polydisc of radius $r'<r$
and bound its integral by that over the four-dimensional box.
This yields
\[
 |h(0)|^2\le(\pi r'^2)^{-2}(2r)^2\,4\exp(1)\|\chi\|_2^2.
\]
Letting $r'\uparrow r$ and taking arbitrary unit vectors proves
\begin{equation}\label{eq:rotation-regularized-point-bound}
 \|\mathcal B_\lambda^\delta(\xi,\xi')\|
       \le\frac{k_2\|\chi\|_2}{d_0-\delta}.
\end{equation}

To identify the limit, put $r_0=d_0/4$.
For $\delta<r_0$ the functions
$h_\delta(z)=\langle\psi,\mathcal B_\lambda^\delta
(\xi+z,\xi'-z)\psi'\rangle$ form a bounded normal family on
$|z^0|,|z^1|<r_0$ by
\eqref{eq:rotation-regularized-point-bound}.
For $\eta\in L^2((-r_0,r_0)^2)$, the real translation identity
and an additional equal time translation on both sides give
\begin{multline}\label{eq:rotation-point-limit-test}
 \int\eta(y)h_\delta(y)\,\dd y\\
 =\langle\psi,T(\xi-2r_0n_0)
       \mathcal B_{2r_0}^{\mathrm c}((\eta*j_\delta)\otimes\lambda)
       T(\xi'-2r_0n_0)\psi'\rangle.
\end{multline}
Since $\eta*j_\delta\to\eta$ in $L^2$, this has a unique limit.
Any two holomorphic subsequential limits thus have the same
integrals on the real square; continuity identifies them there,
and the identity theorem identifies them on the polydisc.
Consequently every matrix element at $z=0$ converges.
The bound \eqref{eq:rotation-regularized-point-bound} defines a
bounded limiting sesquilinear form, represented by
$\mathcal B_\lambda(\xi,\xi')$ with
\eqref{eq:rotation-point-bound}. This is weak operator convergence.
The shift, adjoint, and real linearity identities pass to these
limits.

For holomorphy, fix $(\xi_0,\xi'_0)$ and a positive $\epsilon$
smaller than both tube distances. Near this pair the shift identity
factors the operator as
\[
 T(\xi-\xi_0+\epsilon n_0)\,
 \mathcal B_\lambda(\xi_0-\epsilon n_0,\xi'_0-\epsilon n_0)\,
 T(\xi'-\xi'_0+\epsilon n_0).
\]
The two exterior factors are holomorphic in operator norm.

Finally choose $d'<d$ containing the time support of $\alpha$.
Integrating
\eqref{eq:rotation-point-regularized-translation} with $r=d-\delta$
gives
\[
 \int\alpha(y)\mathcal B_\lambda^\delta(dn_0+y,dn_0-y)\,\dd y
       =\mathcal B_d^{\mathrm c}((\alpha*j_\delta)\otimes\lambda).
\]
The right side converges in norm to the left side of
\eqref{eq:rotation-point-integral}. For matrix elements on the
left, \eqref{eq:rotation-regularized-point-bound} gives the
integrable majorant
$2k_2\|\chi\|_2|\alpha(y)|/(d-d')$ for sufficiently small $\delta$.
Dominated convergence identifies the limiting matrix elements.
Holomorphy makes the limiting integrand norm continuous on the
support, and \eqref{eq:rotation-point-bound} bounds its norm.
It is therefore Bochner integrable, proving the formula and its
stated norm bound.
\end{proof}

\subsection{Ordered field kernels}\label{subsec:rotation-ordered-kernels}

The intermediate arguments of a product of point plates matter
only through their sum. Indeed, for $\kappa_1,\kappa_2\in\mathcal W$
and sufficiently small $\epsilon>0$, the shift identity gives
\[
 \mathcal B_\lambda(\xi,\kappa_1)
       \mathcal B_{\lambda'}(\kappa_2,\xi')
 =\mathcal B_\lambda(\xi,\epsilon n_0)
       T(\kappa_1+\kappa_2-2\epsilon n_0)
       \mathcal B_{\lambda'}(\epsilon n_0,\xi').
\]
The same expression results from any other decomposition of
$\kappa_1+\kappa_2$ into two points of $\mathcal W$.
An exterior argument adjacent to the vacuum can likewise be
changed arbitrarily in $\mathcal W$, since $T(\kappa)\Omega_1=
T(\kappa)^*\Omega_1=\Omega_1$.

For $N\ge2$ and $\xi_1,\ldots,\xi_{N-1}\in\mathcal W$, define
\begin{multline}\label{eq:rotation-ordered-kernel}
 G_N(\xi;\lambda)=
 \langle\Omega_1,\mathcal B_{\lambda_1}(\kappa,\xi_1/2)
       \left(\prod_{a=2}^{N-1}
          \mathcal B_{\lambda_a}(\xi_{a-1}/2,\xi_a/2)\right)\\
       \mathcal B_{\lambda_N}(\xi_{N-1}/2,\kappa)\Omega_1\rangle,
       \qquad\kappa\in\mathcal W.
\end{multline}
Set $G_1(\lambda)=\langle\Omega_1,
\mathcal B_\lambda(\kappa,\kappa)\Omega_1\rangle$.
Both definitions are independent of $\kappa$.
For $x_1,x_{a+1}-x_a\in\mathcal W$, define the vector
\begin{multline}\label{eq:rotation-ordered-vector}
 \Psi(x;\lambda)=
 \mathcal B_{\lambda_1}(x_1,(x_2-x_1)/2)
       \left(\prod_{a=2}^{n-1}
          \mathcal B_{\lambda_a}((x_a-x_{a-1})/2,(x_{a+1}-x_a)/2)\right)\\
       \mathcal B_{\lambda_n}((x_n-x_{n-1})/2,\kappa)\Omega_1.
\end{multline}
Products here and below are ordered with increasing index from
left to right; an empty product is the identity.
For $n=1$ this means
$\mathcal B_{\lambda_1}(x_1,\kappa)\Omega_1$; for $n=0$ it means
$\Omega_1$. It too is independent of $\kappa$.
Put $s(\lambda)=\prod_a s(\sigma_a)$, $w(\lambda)=\sum_a w(\sigma_a)$,
and let $\bar\lambda$ exchange the $+$ and $-$ labels.

\begin{proposition}\label{prop:rotation-ordered-kernels}
The kernels and vectors are holomorphic on their stated domains.
For $N\ge2$,
\begin{equation}\label{eq:rotation-ordered-kernel-bound}
 |G_N(\xi;\lambda)|
 \le\prod_{a=1}^N(2k_2\|\chi_a\|_2)
       \left(\min_a d_{\mathcal W}(\xi_a)\right)^{-N}.
\end{equation}
Let $\alpha_1,\ldots,\alpha_N\in L^2(\R^2)$ have compact
time-ordered supports:
$\sup_{\supp\alpha_a}y^0<\inf_{\supp\alpha_{a+1}}y^0$.
Then
\begin{equation}\label{eq:rotation-ordered-density}
 \E_\mu\prod_a\mathbb B(\alpha_a\otimes\lambda_a)
 =\int\prod_a\alpha_a(y_a)\,
       G_N(y_2-y_1,\ldots,y_N-y_{N-1};\lambda)\,\dd Y.
\end{equation}
If the supports are also in $y^0>0$, their product class is
the Bochner integral of $\prod_a\alpha_a(y_a)\Psi(Y;\lambda)$.
One has $G_1=0$,
$\overline{G_N(\xi;\lambda)}=G_N(\bar\xi;\bar\lambda)$, and
\begin{multline}\label{eq:rotation-ordered-Gram}
 \langle\Psi(x;\lambda),\Psi(z;\lambda')\rangle
 =s(\lambda)G_{n+n'}((x_n-x_{n-1})^*,\ldots,(x_2-x_1)^*,\\
       x_1^*+z_1,z_2-z_1,\ldots,z_{n'}-z_{n'-1};
       \lambda_n,\ldots,\lambda_1,\lambda'_1,\ldots,\lambda'_{n'})
\end{multline}
when $n,n'\ge1$. With a vacuum in the first slot the inner
product is $G_{n'}$ on the successive differences of $z$;
with a vacuum in the second slot it is $s(\lambda)G_n$
on the reversed starred differences and reversed labels.
Furthermore
$T(\kappa')\Psi(x;\lambda)=\Psi(x+\kappa';\lambda)$ for
$\kappa'\in\overline{\mathcal W}$, where every point is translated.
\end{proposition}

\begin{proof}
Holomorphy follows from that of the factors. In the definition of
$G_N$, choose $\kappa=tn_0$ with
$t\ge\frac12\min_a d_{\mathcal W}(\xi_a)$ and apply
\eqref{eq:rotation-point-bound} to every factor. This proves
\eqref{eq:rotation-ordered-kernel-bound}.

To prove the density formula, first put the supports in $y^0>0$.
Choose disjoint time slabs $[a_j,b_j]$, in order, with each support
strictly inside its slab, $a_1\ge0$, and $b_j\le a_{j+1}$.
Let $f_j=\alpha_j\otimes\lambda_j$,
$c_j=(a_j+b_j)/2$, $d_j=(b_j-a_j)/2$, and
$I_j=\mathcal B_{d_j}^{\mathrm c}(\tau_{-c_je_1}f_j)$.
The insertion identity and its localization give
\begin{equation}\label{eq:rotation-ordered-product-class}
 \left[\prod_j\mathbb B(f_j)\right]_1
 =T(a_1n_0)I_1T((a_2-b_1)n_0)I_2\cdots I_N\Omega_1.
\end{equation}
For precision, the backward induction uses
$\Phi'=\prod_{j>j_0}\mathbb B(\tau_{-b_{j_0}e_1}f_j)$,
which belongs to $\mathcal A_+$.
The tail product is $U_{b_{j_0}}\Phi'$, so insertion of $f_{j_0}$
is $\mathcal B_{b_{j_0}}(f_{j_0})[\Phi']_1$.
Localization supplies $T(a_{j_0}n_0)I_{j_0}$; the induction
on the translated tail supplies the next nonnegative gap
$a_{j_0+1}-b_{j_0}$ and the remaining factors.

Use \eqref{eq:rotation-point-integral} to write
\[
 I_j=\int\alpha_j(y_j)
       \mathcal B_{\lambda_j}(y_j-a_jn_0,b_jn_0-y_j)\,\dd y_j.
\]
On the compact supports all these operators are norm bounded;
the $\alpha_j$ are integrable. The product in
\eqref{eq:rotation-ordered-product-class} is therefore one Bochner
integral. Absorb each gap semigroup into a neighbouring point plate.
The inner arguments between plates $j$ and $j+1$ then add to
$y_{j+1}-y_j$. Replace them by equal halves using the product
identity above. The initial argument becomes $y_1$, and the
terminal argument is removed against the vacuum. The integral is
exactly the stated vector formula. Pairing with $\Omega_1$ proves
\eqref{eq:rotation-ordered-density}.
A common real time translation puts any time-ordered supports in
the positive half-plane, without changing their differences or
their moment, proving the general formula.
For $N=1$, this formula and the vanishing mean field give $G_1=0$.

For the Gram identity, take the adjoint of the product defining
$\Psi(x;\lambda)$. It reverses the order, stars the arguments,
and multiplies by $s(\lambda)$. In the middle of the two products
the arguments $x_1^*$ and $z_1$ add to $x_1^*+z_1$; replace them
by equal halves. This gives
\eqref{eq:rotation-ordered-Gram}, including the stated vacuum cases.
The translation identity follows by shifting the first argument.

Conjugating \eqref{eq:rotation-ordered-density} replaces the plates
by $\bar\alpha_a\otimes\bar\lambda_a$. Testing in small ordered
product neighbourhoods identifies
$\overline{G_N(\xi;\lambda)}=G_N(\xi;\bar\lambda)$ on real
$\xi\in\mathcal W^{N-1}$. Both sides, with $\xi$ conjugated on
the first, are holomorphic, and this domain is connected.
The identity theorem proves the claimed conjugation formula.
\end{proof}

\subsection{The two-point angular estimate}
\label{subsec:rotation-two-point-estimate}

The operator bound for $N$ point plates gives angular growth of
order $N$. At two points a positive spectral measure gives the
precise growth order two. Closedness will remove the unwanted
endpoint frequencies.

Fix a coordinate reconstruction $j$. For scalar
$u\in L^1\cap L^2(\R^3)$ supported in $x_j\ge0$, put
\[
 \widetilde u(E,p_\perp)=
       \int u(x)e^{-x_jE+ip_\perp\cdot x_\perp}\,\dd x.
\]
There are complex Radon measures $\varrho^{(j)}_{ab}$,
$a,b\in\{1,2,3\}$, with
\begin{equation}\label{eq:rotation-one-field-spectrum}
 \langle[\mathbb B_a(u)]_j,[\mathbb B_b(u')]_j\rangle
       =\int\overline{\widetilde u}\,\widetilde u'\,
                                    \dd\varrho^{(j)}_{ab}.
\end{equation}
They have a common positive Radon measure $\varrho_0$ such that
$\varrho^{(j)}_{ab}=r_{ab}\varrho_0$ and the matrix $r$ is
positive semidefinite and Hermitian almost everywhere.
In particular, for vector coefficients $c,c'$,
the measures $\varrho_{cc'}=\sum_{a,b}\bar c_ac'_b\varrho_{ab}$
satisfy
\begin{equation}\label{eq:rotation-spectral-CS}
 \int|G_1G_2|\,\dd|\varrho_{cc'}|
 \le\left(\int|G_1|^2\,\dd\varrho_{cc}\right)^{1/2}
       \left(\int|G_2|^2\,\dd\varrho_{c'c'}\right)^{1/2}.
\end{equation}
All these measures are carried by the joint spectral cone.

Here is the construction, including the fact that the measures
need only be locally finite. Convolution and
\exteq{6}{ErgodicComponents} give
\[
 [\mathbb B_a(u*\alpha)]_j
       =\widetilde u(H_j,\mathsf P)[\mathbb B_a(\alpha)]_j
\]
for $\alpha\in L^2$ supported in $x_j\ge0$.
This follows by applying the bounded field and quotient maps to
the $L^2$ Bochner integral
$u*\alpha=\int u(x)\tau_x\alpha\,\dd x$;
Fubini applies to its scalar spectral matrix elements since
$u\in L^1$.
Let $P_1(v)=(\pi(1+v^2))^{-1}$, whose Fourier transform with
either sign is $e^{-|p|}$, and choose
\[
 \alpha_0(x)=\mathbf1_{(0,1)}(x_j)\prod_{i\ne j}P_1(x_i),
 \qquad
 a_0(E,p_\perp)=\widetilde\alpha_0
       =\frac{1-e^{-E}}{E}e^{-|p_\perp|_1}>0.
\]
The quotient at $E=0$ is one. With
$\psi_a=[\mathbb B_a(\alpha_0)]_j$ and
$m_{ab}=\langle\psi_a,\mathsf E_j(\cdot)\psi_b\rangle$,
commutativity of convolution gives
\begin{equation}\label{eq:rotation-spectral-reference}
 a_0(H_j,\mathsf P)[\mathbb B_a(u)]_j
       =\widetilde u(H_j,\mathsf P)\psi_a.
\end{equation}
The finite matrix measure $m$ is positive. If $m_0=\sum_a m_{aa}$,
spectral Cauchy--Schwarz shows $m_{ab}\ll m_0$.
Its density matrix is positive Hermitian on one common full-measure
set: test its quadratic form on the countable rational complex
vectors and then use continuity. Define
$\varrho_{ab}=a_0^{-2}m_{ab}$ and $\varrho_0=a_0^{-2}m_0$.
The strictly positive continuous $a_0$ has a positive lower bound
on bounded sets, so these measures are Radon.
Pointwise matrix Cauchy--Schwarz followed by its integral version
gives \eqref{eq:rotation-spectral-CS}.

On $A_n=\{a_0>1/n\}$, which is bounded, one may divide
\eqref{eq:rotation-spectral-reference} by $a_0$ using bounded
spectral calculus. This proves \eqref{eq:rotation-one-field-spectrum}
with the measures restricted to $A_n$.
As $A_n$ increase to the whole spectral space,
the diagonal identity follows by monotone convergence; its integral
is $\|[\mathbb B_a(u)]_j\|^2\le\|u\|_2^2$.
Equation \eqref{eq:rotation-spectral-CS} then makes the mixed
integral absolutely convergent and gives the unrestricted identity
by dominated convergence. Cone support follows from that of
$\mathsf E_j$.

Now use $j=1$, and take real profiles
$\chi,\chi'\in C_c^\infty(\R)$. Their transverse transforms in this
spectral convention are
$\widehat\chi_+(p_3)=\int\chi(x_3)e^{ip_3x_3}\,\dd x_3$.
For labels $\lambda=(\sigma,\chi)$ and
$\lambda'=(\sigma',\chi')$, set
\[
 \varrho_{\lambda\lambda'}=
 \sum_{a,b}\overline{e^\sigma_a}e^{\sigma'}_b
       \overline{\widehat\chi_+}\widehat{\chi'}_+\,\varrho^{(1)}_{ab}.
\]
In this formula $\widehat{\chi'}_+$ means the same positive-sign
transform of $\chi'$. These are Radon measures on $(E,p_2,p_3)$,
and $\varrho_{\lambda\lambda}\ge0$.
For $\alpha,\alpha'\in L^1\cap L^2(\R^2;\C)$ vanishing almost
everywhere on $y^0<0$, put
$\alpha^L(E,p_2)=\int\alpha(y)e^{-y^0E+iy^1p_2}\,\dd y$.
Equation \eqref{eq:rotation-one-field-spectrum} gives
\begin{equation}\label{eq:rotation-plate-spectrum}
 \langle[\mathbb B(\alpha\otimes\lambda)]_1,
          [\mathbb B(\alpha'\otimes\lambda')]_1\rangle
 =\int\overline{\alpha^L}(\alpha')^L\,
                                     \dd\varrho_{\lambda\lambda'}.
\end{equation}

The crucial estimate is
\begin{equation}\label{eq:rotation-plate-Laplace-bound}
 I_\lambda(t):=\int e^{-tE}\,\dd\varrho_{\lambda\lambda}
       \le\frac{36\|\chi\|_2^2}{\pi t^2}\quad(t>0).
\end{equation}
Indeed, for $c>0$ take
$\alpha_c(y)=c^{-2}\mathbf1_{(c,2c)}(y^0)P_1(y^1/c)$.
On $E\ge|p_2|$ its transform is bounded below by
$e^{-2cE-c|p_2|}\ge e^{-3cE}$.
Since $\|\alpha_c\|_2^2=(2\pi c^2)^{-1}$, the variance bound and
the quotient contraction in \eqref{eq:rotation-plate-spectrum} give
$I_\lambda(6c)\le\|\chi\|_2^2/(\pi c^2)$.
Choosing $c=t/6$ proves the estimate. Its mixed version is
\begin{equation}\label{eq:rotation-plate-mixed-Laplace}
 \int e^{-tE}\,\dd|\varrho_{\lambda\lambda'}|
       \le\sqrt{I_\lambda(t)I_{\lambda'}(t)}.
\end{equation}

Define, for $y^0>0$,
\[
 G^{[1]}_{\lambda\lambda'}(y)
       =s(\sigma)\int e^{-y^0E+iy^1p_2}\,
                                   \dd\varrho_{\lambda\lambda'}.
\]
It is smooth, since on $y^0\ge t>0$ every derivative of order $q$
is bounded in modulus by
$E^qe^{-tE}\le q!(2/t)^qe^{-tE/2}$ on the cone.
It computes the two-point function on time-ordered compact plane
supports. To see this, translate the supports to opposite sides
of $y^0=0$ and put $\alpha^\theta=\bar\alpha\circ r_1$.
Then
\[
 \E_\mu[\mathbb B(\alpha\otimes\lambda)
                 \mathbb B(\alpha'\otimes\lambda')]
 =s(\sigma)\int\overline{(\alpha^\theta)^L}(\alpha')^L\,
                                    \dd\varrho_{\lambda\lambda'}.
\]
Expanding the transforms gives the integral against
$G^{[1]}_{\lambda\lambda'}(y'-y)$. Fubini is justified by
\[
 \|\alpha\|_1\|\alpha'\|_1e^{-\delta E},
\]
where $\delta>0$ is the time gap, and
\eqref{eq:rotation-plate-mixed-Laplace}.
Thus this kernel agrees with $G_2$ on real positive-time differences.

For $r>0$ and $n_\zeta=(\cos\zeta,\sin\zeta)$, it has the
holomorphic continuation
\[
 h^{[1]}_{\lambda\lambda'}(\zeta)
   =s(\sigma)\int
       e^{-rE\cos\zeta+irp_2\sin\zeta}\,
                                \dd\varrho_{\lambda\lambda'}
       \quad(|\Re\zeta|<\pi/2).
\]
If $\zeta=t+i\beta$, the real part of the exponent is
$-r\cos t(E\cosh\beta+p_2\sinh\beta)$.
The cone bounds this above by
$-r\cos t\,e^{-|\beta|}E$.
On every compact substrip this gives an integrable exponential
majorant; dominated convergence and Morera's theorem give
holomorphy. Equations
\eqref{eq:rotation-plate-Laplace-bound}--\eqref{eq:rotation-plate-mixed-Laplace}
give
\begin{equation}\label{eq:rotation-two-point-strip}
 |h^{[1]}_{\lambda\lambda'}(t+i\beta)|
 \le\frac{36\|\chi\|_2\|\chi'\|_2}
          {\pi(r\cos t)^2}\,e^{2|\beta|}.
\end{equation}

\subsection{Closedness removes the unwanted two-point frequencies}
\label{subsec:rotation-two-point-frequencies}

Assume in addition that $\mu$ is closed and cubically invariant.
In particular, its two-point form is invariant under
$R_{\pi/2}$. We use only this quarter-turn consequence here.

For $\upsilon\in(\pi/2)\Z$ and $n_\upsilon\cdot y>0$, define
\[
 G^{(\upsilon)}_{\lambda\lambda'}(y)
   =e^{iw\upsilon}
       G^{[1]}_{\lambda\lambda'}(R_{-\upsilon}y),
       \qquad w=w(\sigma)+w(\sigma').
\]
By quarter-turn invariance and
\eqref{eq:rotation-plate-actions}, this is the two-point density
for supports ordered along $n_\upsilon$.
Two such densities agree on overlaps: test them on small product
neighbourhoods ordered in both directions and use continuity.
The four half-planes cover $\R^2\setminus\{0\}$, so they define
one smooth kernel $G_{\lambda\lambda'}$ there.
It computes the two-point function for any two disjoint compact
plane supports. Indeed, partition them into sufficiently small
half-open squares. Each pair of squares is ordered along one of
the four coordinate directions, and there are only finitely many
pairs. Sum the ordered density identities.

For fixed $r>0$, the functions
\[
 h^{(\upsilon)}(\zeta)=e^{iw\upsilon}
       h^{[1]}_{\lambda\lambda'}(\zeta-\upsilon)
       \quad(|\Re\zeta-\upsilon|<\pi/2)
\]
agree on adjacent strips, by the identity theorem applied on
their real overlap. They give an entire function $h$ with
\begin{equation}\label{eq:rotation-two-point-entire-angle}
 h(\vartheta)=G_{\lambda\lambda'}(rn_\vartheta),\qquad
 h(\zeta+\pi/2)=i^w h(\zeta),\qquad
 |h(t+i\beta)|\le C_{\chi,\chi'}r^{-2}e^{2|\beta|},
\end{equation}
where $C_{\chi,\chi'}=72\|\chi\|_2\|\chi'\|_2/\pi$.
For the bound, choose a strip centre within $\pi/4$ of $t$ and
apply \eqref{eq:rotation-two-point-strip}.

The entire function $e^{-iw\zeta}h(\zeta)$ has period $\pi/2$,
so it is $H(e^{4i\zeta})$ for a holomorphic $H$ on
$\C\setminus\{0\}$. The Laurent series of $H$ gives
frequencies $k=w+4m$, $m\in\Z$. Cauchy's formula on the circle
parametrized by $e^{4i(t+i\beta)}$, $0\le t\le\pi/2$,
bounds the coefficient $a^k(r)$ by
$C_{\chi,\chi'}r^{-2}e^{2|\beta|+k\beta}$.
Sending $\beta$ to the appropriate infinity kills $|k|>2$.
Consequently
\begin{equation}\label{eq:rotation-two-point-window}
 G_{\lambda\lambda'}(rn_\vartheta)
   =\sum_{\substack{|k|\le2\\k-w\in4\Z}}a^k(r)e^{ik\vartheta},
 \qquad |a^k(r)|\le C_{\chi,\chi'}r^{-2}.
\end{equation}
These coefficients are smooth in $r$ and real bilinear in the
profiles, by their Fourier integral formula.

Every pair of labels has only its expected frequency $k=w$,
except for the pairs $(+,+)$ and $(-,-)$, whose windows contain
both $2$ and $-2$. To eliminate the opposite frequency, use
closedness in its vector form $\mathbb B(\nabla u)=0$.
Take $\eta\in C_c^\infty(\R^2;\C)$ and
$\chi_0\in C_c^\infty(\R;\R)$. With $z=y^0+iy^1$ and
$\partial_z=(\partial_{y^0}-i\partial_{y^1})/2$,
\[
 \nabla(\eta(y)\chi_0(x_3))
   =(\partial_z\eta)\otimes(+,\chi_0)
       +(\partial_{\bar z}\eta)\otimes(-,\chi_0)
       +\eta\otimes(3,\chi'_0).
\]
Pair this gradient with a second plate whose plane support is
disjoint. Apply the density identities and integrate by parts.
Testing on small disjoint product neighbourhoods gives
\begin{equation}\label{eq:rotation-two-point-Bianchi}
 \partial_zG_{(+,\chi_0)\lambda'}
       +\partial_{\bar z}G_{(-,\chi_0)\lambda'}
       +G_{(3,\chi'_0)\lambda'}=0
                       \quad\text{on }\R^2\setminus\{0\}.
\end{equation}
There is no diagonal contribution in this equation.
In polar coordinates,
\begin{align*}
 \partial_z(f(r)e^{ik\vartheta})
       &=\tfrac12e^{i(k-1)\vartheta}(f'(r)+kf(r)/r),\\
 \partial_{\bar z}(f(r)e^{ik\vartheta})
       &=\tfrac12e^{i(k+1)\vartheta}(f'(r)-kf(r)/r).
\end{align*}
For $\lambda'=(+,\chi')$, the frequency $-3$ in
\eqref{eq:rotation-two-point-Bianchi} occurs only in the derivative
of $a^{-2}_{(+,\chi_0)(+,\chi')}$.
It obeys $f'=2f/r$.
For $\lambda'=(-,\chi')$, the frequency $3$ likewise gives
$f'=2f/r$ for $a^2_{(-,\chi_0)(-,\chi')}$.
Each solution is $f(r)=f(1)r^2$. The bound
$|f(r)|\le C_{\chi_0,\chi'}r^{-2}$ for all $r>0$ forces
$f(1)=0$ by sending $r\to\infty$. Thus
\begin{equation}\label{eq:rotation-two-point-phase}
 G_{\lambda\lambda'}(R_ty)
       =e^{i(w(\sigma)+w(\sigma'))t}G_{\lambda\lambda'}(y).
\end{equation}
The plate action cancels this phase in the two-point density
formula, proving rotational covariance for disjoint compact
plane supports.

\subsection{Extending identities across coincident plane supports}
\label{subsec:rotation-contact}

The preceding density argument does not yet address coincident
plane supports. We use the following $L^2$ argument,
which will also be applied at every higher order.

\begin{proposition}\label{prop:rotation-no-contact}
Let $U\subset\R^2$ be open, and write $L^2_U$ for the vector
fields supported in $U\times\R$. Let $D:(L^2_U)^N\to\C$ be a
bounded multilinear form. Suppose it vanishes when the entries
are supported over pairwise disjoint compact subsets of $U$,
and whenever one entry is $\nabla u$ with
$u\in C_c^\infty(U\times\R)$. Then $D=0$.
\end{proposition}

\begin{proof}
We argue simultaneously for all open $U$ by induction on $N$.
At $N=1$, fields with compact plane support are dense, so the
first assumption suffices.

At $N=2$, represent the bilinear form as
$D(f_1,f_2)=\int f_1(x)\cdot(Tf_2)(x)\,\dd x$ for a bounded
operator $T$ on $L^2_U$. The pairing here is bilinear.
For disjoint compact $K,K'\subset U$ the first assumption gives
$\mathbf1_KT\mathbf1_{K'}=0$.
Inner regularity and strong convergence of the corresponding
multiplication projections extend this to disjoint Borel sets.
Using $A$ and $U\setminus A$ on both sides gives
$T\mathbf1_A=\mathbf1_AT$ for every Borel $A$.
Uniform approximation by simple functions then gives $Th=hT$
for every bounded plane multiplier $h$. Hence
\[
 D(h\nabla u,f_2)=D(\nabla u,hf_2)=0.
\]

Fix a compact $A\subset U$, choose $u_0\in C_c^\infty(U)$ equal
to one near $A$, and let $\chi_0\in C_c^\infty(\R)$.
On $A\times\R$, the gradients of
$u_0\chi_0$, $u_0y^0\chi_0$, and $u_0y^1\chi_0$ are respectively
\[
 \chi_0'e_3,\qquad
 \chi_0e_1+y^0\chi_0'e_3,\qquad
 \chi_0e_2+y^1\chi_0'e_3.
\]
The preceding identity with multiplier $h\mathbf1_A$ first
annihilates $h\mathbf1_A\chi'_0e_3$.
Applying it with $hy^0\mathbf1_A$ and $hy^1\mathbf1_A$ then
annihilates $h\mathbf1_A\chi_0e_1$ and
$h\mathbf1_A\chi_0e_2$. These multipliers are bounded since $A$
is compact. Such plane-transverse products span a dense subspace.
For the transverse derivative family, its orthogonal complement
in $L^2(\R)$ has distributional derivative zero, hence is a
constant and therefore zero. Thus these three vector families
are total in $L^2_U$, and $D=0$.

Suppose $N\ge3$ and the assertion is proved at every smaller order.
Two consequences of the induction hypothesis separate unequal
plane cells. First, if one entry is supported over a compact
$K\subset U$ and the others over $U\setminus K$, fix that entry.
The form on the remaining $N-1$ slots satisfies both assumptions
on the open set $U\setminus K$ and is zero.
Second, if the slots divide into nonempty groups $I,J$, with
the $J$ entries supported over a compact $K_J$ and the $I$
entries over its complement, fix the $J$ entries.
The form on the $I$ slots inherits the gradient assumption.
When those slots have pairwise disjoint compact supports, the
first consequence, applied to any one of them, gives zero.
The induction hypothesis on $U\setminus K_J$ therefore makes
this form zero as well.

By density put all entries over one compact $K_0\subset U$.
Partition the plane into half-open squares $Q$ of side $d$ and
expand $D$ by multilinearity. A tuple of squares which are not
all equal has zero contribution. To justify compact separation,
replace each square by its concentric closed square of side
$d-2\epsilon$. Separate the slots whose square equals the first
from the remaining slots, use the second consequence with
$K_J=K_0$ intersected with the union of their closed squares,
and send $\epsilon\downarrow0$ in $L^2$.
Only the equal-square terms remain, and Cauchy--Schwarz gives
\[
 |D(f_1,\ldots,f_N)|
 \le\|D\|\max_Q\|\mathbf1_Qf_3\|_2
       \|f_1\|_2\|f_2\|_2\prod_{i=4}^N\|f_i\|_2.
\]
Indeed, sum the first two square norms by Cauchy--Schwarz and
bound the others by their whole norms. The maximum tends to
zero as $d\downarrow0$: its square is the integral over an
area-$d^2$ cell of the fixed integrable function
$y\mapsto\int|f_3(y,x_3)|^2\,\dd x_3$.
Absolute continuity of the integral proves the assertion.
\end{proof}

Apply this proposition with $N=2$ to
\[
 D(b,b')=\E_\mu[\mathbb B(R_tb)\mathbb B(R_tb')]
              -\E_\mu[\mathbb B(b)\mathbb B(b')],
\]
where $R_tb$ means $R_t\cdot b$. Gaussian domination makes it a
bounded bilinear form. Finite sums of plates with smooth real
transverse profiles are dense in the fields over any fixed
compact plane set: simple plane-transverse tensors are dense
in $L^2$, and $e^+,e^-,e^3$ span $\C^3$.
Thus \eqref{eq:rotation-two-point-phase} supplies the disjoint-support
assumption for general fields. Closedness supplies the gradient
assumption, since rotations carry gradients to gradients.
Consequently the two-point function is invariant under every
$R_t$, on all $L^2$ tests.

\subsection{The angular expansion at every order}
\label{subsec:rotation-angular-expansion}

Fix $N\ge2$ and a tuple of labels $\lambda$, with total weight
$w=w(\lambda)$. Put $\mathcal W_\upsilon=R_\upsilon\mathcal W$
for $\upsilon\in(\pi/2)\Z$. For a permutation $\mathfrak s$,
let
\[
 \mathcal Z(\mathfrak s,\upsilon)=
 \{Z=(z_1,\ldots,z_N):
       z_{\mathfrak s(a+1)}-z_{\mathfrak s(a)}
                    \in\mathcal W_\upsilon,\ 1\le a<N\}.
\]
On this open convex set define
\[
 G_\upsilon^{\mathfrak s}(Z;\lambda)=e^{iw\upsilon}
 G_N\left(
       (R_{-\upsilon}(z_{\mathfrak s(a+1)}-z_{\mathfrak s(a)}))_{a=1}^{N-1};
       \mathfrak s\lambda\right).
\]
Here $(\mathfrak s\lambda)_a=\lambda_{\mathfrak s(a)}$.
Moment symmetry, quarter-turn invariance, and
\eqref{eq:rotation-ordered-density} identify this as the moment
density on compact product supports in its real part.
Two such holomorphic functions agree on their intersection.
Indeed, the intersection is convex and is preserved by taking
real parts, so a nonempty intersection contains an open set of
real configurations. Testing both densities there identifies
them by continuity, and the holomorphic identity theorem
identifies them on the full intersection.
They therefore define one holomorphic function
$\mathcal G_N(Z;\lambda)$ on
$\mathcal Z_N=\bigcup_{\mathfrak s,\upsilon}\mathcal Z(\mathfrak s,\upsilon)$.

The real set $\mathcal O_N=\mathcal Z_N\cap(\R^2)^N$ consists
of configurations whose first coordinates are all distinct
or whose second coordinates are all distinct.
Its complement is the finite union of the codimension-two
subspaces
$\{y_a^0=y_b^0,\ y_k^1=y_l^1\}$.
Thus $\mathcal O_N$ is open, dense, and connected.
The functions $\mathcal G_N$ are symmetric under simultaneous
permutation of points and labels, invariant under common
translations, and satisfy
\[
 \mathcal G_N(R_{\pi/2}Z;\lambda)=i^w\mathcal G_N(Z;\lambda),
 \qquad
 \overline{\mathcal G_N(Z;\lambda)}
       =\mathcal G_N(\bar Z;\bar\lambda).
\]
These identities follow from the defining densities and
the corresponding identities for $G_N$.

For $\beta\in\R$ write $\Lambda_\beta=R_{i\beta}$.
In the coordinates $\xi^\pm=\xi^0\pm i\xi^1$,
$\mathcal W=\{\Re\xi^+>0,\Re\xi^->0\}$ and
$\Lambda_\beta\xi^\pm=e^{\mp\beta}\xi^\pm$.
Thus every tube $\mathcal Z(\mathfrak s,\upsilon)$ is preserved
by $\Lambda_\beta$. Moreover
$d_{\mathcal W}(\Lambda_\beta\xi)\ge
e^{-|\beta|}d_{\mathcal W}(\xi)$ on the closed tube.

\begin{proposition}\label{prop:rotation-finite-angular-expansion}
There are holomorphic coefficients $c_m(Z;\lambda)$ on
$\mathcal Z_N$, vanishing unless $|w+4m|\le N$, such that
\begin{equation}\label{eq:rotation-finite-angular-expansion}
 \mathcal G_N(\Lambda_\beta Z;\lambda)
       =\sum_m c_m(Z;\lambda)e^{-(w+4m)\beta}
                         \quad(\beta\in\R).
\end{equation}
They are unique, symmetric under simultaneous permutation,
and invariant under common translations. They satisfy
\begin{align}
 c_m(\Lambda_\beta Z;\lambda)
       &=e^{-(w+4m)\beta}c_m(Z;\lambda),
                           \label{eq:rotation-angular-coefficient-identities}\\
 c_m(R_{\pi/2}Z;\lambda)&=i^wc_m(Z;\lambda),\qquad
 \sum_m c_m(Z;\lambda)=\mathcal G_N(Z;\lambda),\notag\\
 \overline{c_m(X;\lambda)}&=c_{-m}(X;\bar\lambda)
                           \quad(X\in\mathcal O_N).\notag
\end{align}
\end{proposition}

\begin{proof}
Fix $0<\vartheta_0<\pi/4$ and first consider real configurations
whose consecutive differences have positive lengths and angles
$|\vartheta_a|<\vartheta_0$ from $n_0$.
For $\zeta=t+i\beta$,
\[
 \Re(R_{\zeta-\upsilon}\xi_a)^\pm
       =|\xi_a|e^{\mp\beta}\cos(\vartheta_a+t-\upsilon).
\]
The strips $|t-\upsilon|<\pi/2-\vartheta_0$ cover the complex
$\zeta$-plane. Thus $h_X(\zeta)=\mathcal G_N(R_\zeta X;\lambda)$
is entire. Choosing $\upsilon$ within $\pi/4$ of $t$ bounds every
gap distance below by
$r_Xe^{-|\beta|}\cos(\vartheta_0+\pi/4)$, where
$r_X=\min_a|\xi_a|$.
Equation \eqref{eq:rotation-ordered-kernel-bound} gives
$|h_X(t+i\beta)|\le C_{X,\lambda}e^{N|\beta|}$.
Its quarter-period is $h_X(\zeta+\pi/2)=i^wh_X(\zeta)$.
The Laurent coefficient argument used at two points, with
growth exponent $N$, gives the finite frequencies
$k=w+4m$, $|k|\le N$.

To extend this statement to a whole tube, let
$\mathcal M=\{m:|w+4m|\le N\}$ and
$P(u)=\prod_{m\in\mathcal M}(u+w+4m)$.
The functions
$Z\mapsto P(\partial_\beta)G_\upsilon^{\mathfrak s}
(\Lambda_\beta Z;\lambda)$ are holomorphic for each real $\beta$.
They vanish on the nonempty open real subset obtained by rotating
and permuting the configurations just used.
The identity theorem gives zero on the entire convex tube.
The resulting constant-coefficient ordinary differential equation
has distinct characteristic roots $-(w+4m)$.
Its coefficients in the exponential expansion are fixed linear
combinations of its derivatives at $\beta=0$, and hence are
holomorphic in $Z$.
On overlapping tubes the exponentials are linearly independent,
so the coefficients agree and define $c_m$ on $\mathcal Z_N$.
This also proves uniqueness. Comparing coefficients after a
permutation, common translation, or addition of the parameter
$\beta$ gives the stated identities. For the conjugation identity
use $\overline{\Lambda_\beta X}=\Lambda_{-\beta}X$ and
$w(\bar\lambda)=-w$.
\end{proof}

The coefficients have exact real rotation phases whenever both
configurations are in $\mathcal O_N$:
\begin{equation}\label{eq:rotation-angular-coefficient-phase}
 c_m(R_tX;\lambda)=e^{i(w+4m)t}c_m(X;\lambda).
\end{equation}
Locally in one tube, the allowed angles form an interval:
each ordered difference has a positive projection precisely
when $\cos(\vartheta_a+t)>0$.
The identity theorem in $t+i\beta$ extends
\eqref{eq:rotation-finite-angular-expansion} from the imaginary
axis to that interval. Comparing the exponential coefficients
gives \eqref{eq:rotation-angular-coefficient-phase} locally.

To pass through changes of ordering, first take distinct points
with no two distinct chords perpendicular. Such configurations
are a dense open set, the complement of finitely many nonzero
polynomial zero sets. Every rotation of one of them belongs to
$\mathcal O_N$: simultaneous coincidences of first and second
coordinates would give perpendicular chords, or identical points.
The local phase identity therefore continues along its whole
rotation orbit. Approximation proves it whenever $X,R_tX$
are in $\mathcal O_N$.

For an arbitrary configuration $X$ of distinct plane points,
choose $t$ so that all first coordinates of $R_tX$ are distinct.
Only finitely many angles modulo $\pi$ are excluded.
Define
\[
 \widetilde c_m(X;\lambda)
       =e^{-i(w+4m)t}c_m(R_tX;\lambda).
\]
The phase identity makes this independent of the choice.
Locally one may keep $t$ fixed, so it is real analytic on the
connected space $(\R^2)^N_{\ne}$ of distinct configurations.
It inherits permutation symmetry, translation invariance, and
\[
 \widetilde c_m(R_tX;\lambda)
       =e^{i(w+4m)t}\widetilde c_m(X;\lambda).
\]
Put $\widetilde G_N=\sum_m\widetilde c_m$.
It is real analytic and agrees with $\mathcal G_N$ on $\mathcal O_N$.

Define $\mathfrak a(0)=\mathfrak a(1)=0$, and, for $N\ge2$,
let $\mathfrak a(N)$ be the largest $|m|$ for which $c_m$ is
not identically zero for some tuple of $N$ labels; set it to zero
if all coefficients vanish. Since $|w|\le N$ and $|w+4m|\le N$,
\begin{equation}\label{eq:rotation-angular-linear-bound}
 \mathfrak a(N)\le N/2.
\end{equation}
A nonzero coefficient can be detected on a strictly time-ordered
real configuration after a quarter turn and permutation:
the identity theorem on each tube and the coefficient identities
above give this reduction. Conjugation then shows that if
$\mathfrak a(N)>0$, both indices
$\mathfrak a(N)$ and $-\mathfrak a(N)$ occur.

\subsection{Reflection bounds for the angular coefficients}
\label{subsec:rotation-angular-reflections}

We next prove the convexity estimate
\begin{equation}\label{eq:rotation-angular-convexity}
 2\mathfrak a(a+b)\le\mathfrak a(2a)+\mathfrak a(2b)
                          \quad(a,b\ge0).
\end{equation}
Let $N=a+b$ and $M=\mathfrak a(N)>0$.
Choose a strictly time-ordered $X=(y_1,\ldots,y_N)$ and labels
with $c_{-M}(X;\lambda)\ne0$. Translate its time cut so that the
first $a$ points have negative times and the remaining $b$
have positive times. Reflect the first group and reverse its
order to obtain $x$, with reversed labels $\lambda^x$;
call the second group $z$, with labels $\lambda^z$.
Their weights satisfy $w=w_x+w_z$.
The Gram identity gives
\[
 \langle\Psi(\Lambda_\beta x;\lambda^x),
             \Psi(\Lambda_\beta z;\lambda^z)\rangle
       =s(\lambda^x)\mathcal G_N(\Lambda_\beta X;\lambda).
\]
An empty group gives the vacuum.
The squared norm of the first vector is, up to the sign
$s(\lambda^x)$, the kernel of its reflected double:
the configuration $(r_1x_a,\ldots,r_1x_1,x_1,\ldots,x_a)$
with the same labels in reverse then forward order.
Its weight is $2w_x$. The other squared norm has weight $2w_z$.
These are direct instances of
\eqref{eq:rotation-ordered-Gram}; the reflection-conjugation
keeps the labels.

Apply Cauchy--Schwarz to this matrix element and expand the
three kernels by \eqref{eq:rotation-finite-angular-expansion}.
Multiplication by $e^{2w\beta}$ cancels their weights and gives
\[
 \left|\sum_m c_m(X;\lambda)e^{-4m\beta}\right|^2
 \le
 \left|\sum_{|m|\le\mathfrak a(2a)}d_me^{-4m\beta}\right|
 \left|\sum_{|m|\le\mathfrak a(2b)}d'_me^{-4m\beta}\right|.
\]
The coefficients $d_m,d'_m$ belong to the two specified doubles.
A vacuum factor is one. As $\beta\to+\infty$, the left side
has a nonzero leading term of rate $8M$, and the right side is
bounded by a constant times
$e^{4(\mathfrak a(2a)+\mathfrak a(2b))\beta}$.
This proves \eqref{eq:rotation-angular-convexity}.
The cases $M=0$ or $N\le1$ already satisfy it.

There is also a positive kernel for the highest coefficient.
For $a\ge1$, set $M=\mathfrak a(2a)$. On labelled configurations
of $a$ distinct points in $y^0>0$, define
\begin{equation}\label{eq:rotation-top-coefficient-kernel}
 K_a((x;\lambda),(z;\lambda'))=
       s(\lambda)\widetilde c_{-M}^{(2a)}
                        (r_1x,z;\lambda,\lambda').
\end{equation}
The superscript indicates the number of points.
This kernel is Hermitian and positive semidefinite.
First suppose the times in each configuration are distinct and
order them increasingly. By the Gram formula,
\[
 \langle\Psi(\Lambda_\beta x;\lambda),
                 \Psi(\Lambda_\beta z;\lambda')\rangle
 =s(\lambda)\sum_{|m|\le M}
       \widetilde c_m^{(2a)}(r_1x,z;\lambda,\lambda')
       e^{-(w(\lambda)+w(\lambda')+4m)\beta}.
\]
Reversing the reflected points is harmless by labelled symmetry.
Multiply this Gram kernel by the positive factors
$e^{w(\lambda)\beta}e^{w(\lambda')\beta}e^{-4M\beta}$.
Its pointwise limit as $\beta\to+\infty$ is $K_a$.
Positive semidefiniteness passes to this limit on every finite
matrix. Approximating configurations with tied times by ones
with distinct times proves the same assertion for all the
stated configurations.

The exact phase of $\widetilde c_{-M}$ transports this consequence
to any line $L$ in the plane. If $x,z$ lie on one open side of
$L$ and $r_L$ is its reflection, translation to a point of $L$
and rotation of its normal to $n_0$ express
$\widetilde c_{-M}(r_Lx,z)$ as $K_a$ times a unit-modulus factor.
Taking absolute values of the positive-kernel
Cauchy--Schwarz inequality gives
\begin{multline}\label{eq:rotation-top-coefficient-CS}
 |\widetilde c_{-M}^{(2a)}(r_Lx,z;\lambda,\lambda')|^2\\
 \le
 |\widetilde c_{-M}^{(2a)}(r_Lx,x;\lambda,\lambda)|
 |\widetilde c_{-M}^{(2a)}(r_Lz,z;\lambda',\lambda')|.
\end{multline}
In particular a zero reflected diagonal forces every such
off-diagonal coefficient to vanish.

\subsection{Regular polygons and the reduction to two points}
\label{subsec:rotation-polygons}

We prove the arithmetic restriction
\begin{equation}\label{eq:rotation-polygon-divisibility}
 a\mid 2\mathfrak a(2a)\quad(a\ge1).
\end{equation}
Fix $a$ and put $M=\mathfrak a(2a)$,
$c_*=\widetilde c_{-M}^{(2a)}$.
The coefficients depend only on the labelled set of points.
If a rotation through $t$ permutes such a set without changing
its labels, its exact phase gives
$(1-e^{i(w-4M)t})c_*=0$.
For a regular $2a$-gon with the same label $(\sigma,\chi)$ at
each vertex, use $t=\pi/a$ and $w=2aw(\sigma)$.
The label weight cancels, leaving $e^{-4\pi iM/a}$.
Thus $c_*$ vanishes on this polygon whenever $a$ fails to divide $2M$.

That selection rule concerns a special configuration. The
positive kernel propagates its zero to arbitrary configurations
by the following halving argument. For a labelled set $S$ of
at most $2a$ distinct points, let $\mathcal N(S)$ mean that
$c_*$ vanishes on every $2a$-point labelled configuration
containing $S$.
Let $S'$ have $j'\le a$ points on one open side $L_+$ of a line
$L$, and suppose $S\subset r_LS'\cup S'$.
Then
\begin{equation}\label{eq:rotation-halving}
 \mathcal N(S)\ \Longrightarrow\
       \mathcal N(S')\ \text{and}\ \mathcal N(r_LS').
\end{equation}
To see this, add $a-j'$ arbitrary points $R$ in $L_+$ to $S'$,
forming $x=S'\cup R$. The reflected double $r_Lx\cup x$
contains $S$, so its diagonal coefficient is zero.
Equation \eqref{eq:rotation-top-coefficient-CS} implies
$c_*(r_Lx\cup z)=0$ for every $a$-point labelled set $z$ in $L_+$.
Thus, with $r_LS'$ fixed, the remaining $a-j'$ points on
$L_-$ and $a$ points on $L_+$ may vary in an open set.

For each choice of their labels, the space of their distinct
positions, avoiding the fixed points, is the complement of
finitely many codimension-two affine subspaces and is connected.
The coefficient is real analytic there. Its zero on the
nonempty open set just obtained extends to all positions.
Every $2a$-point configuration containing $r_LS'$ is obtained
this way. This proves $\mathcal N(r_LS')$.
Interchange the sides of $L$ and repeat the argument to obtain
$\mathcal N(S')$.

Suppose now $a\nmid2M$. Then $M>0$ and $a\ge3$.
Fix an arbitrary label and any regular $2a$-gon with vertices
$P_k$, $k\in\Z/(2a)\Z$, in cyclic order.
Let $\mathcal N_j$ assert $\mathcal N(S)$ for every arc of $j$
consecutive vertices, with this label, in every such polygon.
Selection gives $\mathcal N_{2a}$.
For $j\ge2$, put $j'=\lceil j/2\rceil\le a$ and take an arc
$S'=\{P_{k_0+1},\ldots,P_{k_0+j'}\}$.
Reflect in the perpendicular bisector of
$[P_{k_0},P_{k_0+1}]$.
The $a$ vertices $P_{k_0+1},\ldots,P_{k_0+a}$ lie on one
open side: their angles from the edge-midpoint radius are
$(2k-1)\pi/(2a)\in(0,\pi)$.
Their images satisfy $r_LP_{k_0+k}=P_{k_0+1-k}$.
Consequently $r_LS'\cup S'$ is an arc of $2j'$ consecutive
vertices and contains a $j$-vertex arc.
Equation \eqref{eq:rotation-halving} gives
$\mathcal N_j\Rightarrow\mathcal N_{j'}$.

Iteration reaches $j=1$. Any plane point is a vertex of some
regular $2a$-gon, so the coefficient vanishes whenever a
configuration has even one point with the fixed label.
Since the label was arbitrary, $c_*$ vanishes identically for
all label tuples. This contradicts the occurrence of the
index $-M$, proving \eqref{eq:rotation-polygon-divisibility}.

Now set $\nu(a)=\mathfrak a(2a)$. It is a nonnegative integer
sequence, with $\nu(0)=0$, $\nu(a)\le a$, and
$2\nu(a)\le\nu(a-1)+\nu(a+1)$ by
\eqref{eq:rotation-angular-convexity}.
Its increments $d_a=\nu(a)-\nu(a-1)$ are nondecreasing integers.
The linear bound keeps them bounded above, so they are eventually
constant, say $d$. Thus $\nu(a)=da-c$ for large $a$, where
$c=\sum_k(d-d_k)\ge0$ is an integer.
Divisibility \eqref{eq:rotation-polygon-divisibility} makes every
large $a$ divide $2c$, forcing $c=0$.
All increments equal $d=\nu(1)$, and
\[
 \mathfrak a(2a)=a\mathfrak a(2),\qquad
 2\mathfrak a(N)\le N\mathfrak a(2).
\]
For the second inequality split $N=\lfloor N/2\rfloor+
\lceil N/2\rceil$ in \eqref{eq:rotation-angular-convexity}.

The two-point result gives $\mathfrak a(2)=0$.
For smooth transverse profiles this follows from
\eqref{eq:rotation-two-point-phase} and the ordered density formula:
near any strictly ordered two-point configuration, only the
frequency $w$ occurs. Linear independence of the finite angular
exponentials makes $c_m=0$ there for $m\ne0$, and the identity
theorem on each tube gives zero throughout.
For arbitrary $L^2$ profiles, the all-$L^2$ two-point covariance
proved above gives the same local density identity.
Hence
\begin{equation}\label{eq:rotation-no-extra-angular-coefficients}
 \mathfrak a(N)=0\quad\text{for every }N,\qquad
 \widetilde G_N(R_tX;\lambda)=e^{iw(\lambda)t}
                                  \widetilde G_N(X;\lambda).
\end{equation}

\subsection{Moment densities on all distinct configurations}
\label{subsec:rotation-all-distinct-density}

The analytic continuation $\widetilde G_N$ must still be
identified as the moment density at distinct configurations
outside $\mathcal O_N$. This is a separate step; continuation
of a density does not by itself exclude an additional singular
part on that smaller set.

We first record the analytic displacement used for this step.
If polynomial field observables $\Phi'$ and $\Phi''$ have test
supports in $x_1<c-\delta$ and $x_1>c+\delta$, respectively,
then
\[
 s\longmapsto\E_\mu[\Phi'U(se_2)\Phi'']
\]
extends holomorphically to $|\Im s|<2\delta$.
Translate $c$ to zero and put
$\psi'=[U_{-\delta}\Theta_1\Phi']_1$,
$\psi''=[U_{-\delta}\Phi'']_1$.
The quotient and semigroup identities express the real
correlation as
$\langle\psi',T((2\delta,s))\psi''\rangle$.
This is holomorphic on the stated strip by the cone.

For compact, pairwise disjoint plane supports we claim
\begin{equation}\label{eq:rotation-distinct-density}
 \E_\mu\prod_a\mathbb B(\alpha_a\otimes\lambda_a)
 =\int\prod_a\alpha_a(y_a)\widetilde G_N(Y;\lambda)\,\dd Y.
\end{equation}
The integral is absolutely convergent: the product support
is a compact subset of $(\R^2)^N_{\ne}$, the kernel is bounded
there, and each $\alpha_a\in L^2$ is integrable on its compact
support.
Call the difference of the two sides $D(\alpha_1,\ldots,\alpha_N)$
on these tuples.

We show local vanishing near each distinct configuration $X$
by induction on the number of equal-second-coordinate pairs.
When there are none, $X\in\mathcal O_N$ and the ordered tube
density gives the identity in small product discs.
Otherwise take a pair with equal second coordinates and
different times. Choose a time cut between them, different
from all point times, and split the points into groups $I,J$
below and above the cut.
Move the $J$ group by $se_2$. The points remain distinct for
every real $s$, since the groups are separated in time.
A sufficiently small positive $s_0$ removes every
second-coordinate coincidence across the cut and creates none:
take it smaller than every nonzero initial difference of
second coordinates.
There are fewer coincident pairs at $X(s_0)$.

The induction hypothesis gives a radius $2r$ near $X(s_0)$
on which $D$ vanishes. Reduce $r$ so that the original discs
of radius $r$ are disjoint and preserve the time gap.
Fix tests in these original discs, and translate those in
group $J$. The resulting difference $F(s)$ is real analytic
on $\R$.
For its moment term this follows from the separated-cluster
strip argument above. For its density term, change variables
to integrate
$\widetilde G_N(Y+sv_J;\lambda)$ against the fixed product
of tests, where $(v_J)_a=e_2$ for $a\in J$ and zero otherwise.
Over any bounded real $s$-interval these configurations form
a compact set of distinct points. Finite local holomorphic
extensions of the real analytic kernel give a common radius
$\epsilon>0$ and a bound $C$ there. Cauchy's estimates bound
the integrated derivative of order $q$ by
$Cq!\epsilon^{-q}\prod_a\|\alpha_a\|_1$.
Taylor's formula therefore gives real analyticity of this
integral.

For $|s-s_0|<r$ the translated tests remain in the discs of
radius $2r$ at $X(s_0)$, so $F(s)=0$.
The real analytic identity theorem gives $F(0)=0$.
This proves local vanishing near $X$ and completes the induction.
For general disjoint compact supports, finitely many of these
product-disc neighbourhoods cover their product.
Partition each support into sufficiently small half-open
squares. Every product of nonzero pieces lies in one of the
larger vanishing neighbourhoods: choose a point of the
product support in it and make each square diameter smaller
than the available half-radius.
Multilinearity and finite summation prove
\eqref{eq:rotation-distinct-density}.

\subsection{Euclidean covariance and all reflection planes}
\label{subsec:rotation-conclusion}

\begin{theorem}\label{thm:rotation-covariance}
Let $\mu$ be a Gaussian-dominated law of closed tempered two-forms.
Suppose it is invariant under translations and the cubic group,
and reflection positive at the nine planes
$x_j=0$ and $x_j\pm x_i=0$, $i<j$.
Then it is invariant under all Euclidean motions and reflection
positive at every plane, on the full $L^2$ space of either open side.
\end{theorem}

\begin{proof}
Proposition~\extref{3.1}{ErgodicComponents} supplies the sharp coordinate
cones. For the plane $(1,2)$, fix $t$ and define the bounded
$N$-linear form
\[
 D(b_1,\ldots,b_N)=
 \E_\mu\prod_a\mathbb B(R_t\cdot b_a)
       -\E_\mu\prod_a\mathbb B(b_a).
\]
Its boundedness follows from Proposition~\extref{2.1}{ContinuumLimits}.
For plates over pairwise disjoint compact plane supports,
\eqref{eq:rotation-distinct-density} and
\eqref{eq:rotation-no-extra-angular-coefficients} prove $D=0$:
the factor $e^{-iw(\lambda)t}$ from the plate action cancels
the factor $e^{iw(\lambda)t}$ from the kernel.
Density of finite sums of plates gives the same assertion for
arbitrary $L^2$ fields over those compact sets.
The form also vanishes on gradients, by closedness and the
fact that rotations carry gradients to gradients.
\Cref{prop:rotation-no-contact} therefore gives $D=0$
on all $L^2$ entries, for every $N$.

The entire generating function of each finite real Schwartz
test family is determined by these moments. Thus its
characteristic function is unchanged by $R_t$, and Fourier
uniqueness on the evaluation sigma algebra gives
$(R_t)_*\mu=\mu$.
Permuting coordinates applies the same argument in the other
coordinate planes: the nine mirrors give their cones, and
cubic invariance gives their quarter turns.
These plane rotations generate $SO(3)$.
A coordinate reflection already preserves $\mu$, so the full
$O(3)$ acts invariantly. Translation invariance gives the
Euclidean group.

Let $g$ carry $x_1=0$ to any chosen plane and its positive side
to the chosen open side. The composition operator of $g$ is
unitary on $L^2(\mu)$, maps these side spaces onto each other,
and conjugates the axis reflection to the chosen reflection.
Transporting its nonnegative reflection form proves the final
assertion.
\end{proof}

\Cref{thm:rotation-covariance} applies separately to the
component laws of Sections~\extref{5--6}{ErgodicComponents}.
They retain Gaussian domination, closedness, the cubic group,
and the nine mirror forms. Together with their all-direction
clustering, it supplies the hypotheses for the tensor OS
verification.



